\section{LLaVA：Vision Encoder + Projector + Language Model}

本节从 representation learning 进入生成式问答。LLaVA 冻结或初始化预训练 vision encoder 与 LLM，用一个 projector 把视觉 patch embeddings 映射到 LM hidden dimension，再把它们作为一段 visual tokens 拼入 prompt。

\subsection{Architecture 与 shape flow}

设 vision encoder 输出 $V\in\mathbb{R}^{N_v\times d_v}$，projector $W\in\mathbb{R}^{d_v\times d_{\mathrm{lm}}}$，则 visual tokens 为

$$
H_v=VW\in\mathbb{R}^{N_v\times d_{\mathrm{lm}}}.
$$

随后与 text embeddings $H_t$ 在 sequence axis 拼接：$H=[H_v;H_t]$。这里没有把图像“翻译成文字”，而是让 LM attention 直接读取连续视觉 tokens。

\begin{figure}[H]
\centering
\includegraphics[width=0.84\textwidth]{llava-architecture.png}
\caption{LLaVA architecture：CLIP vision encoder、projection layer 与 autoregressive LM。}
\figsource{Stanford CS336 2026 Lecture 17 官方可执行讲义，\texttt{images/llava-architecture.png}。}
\end{figure}

\begin{knowledgebox}{读图：projector 的职责}
Projector 只负责 dimension/interface alignment，不自动保证语义对齐。训练第一阶段通常冻结两端，先让 projector 学会把视觉 features 放到 LM 可读取的空间；第二阶段再用 instruction data 调整 LM 或全模型。
\end{knowledgebox}

\subsection{Training stages 与 synthetic instruction data}

这里从接口转向数据。LLaVA 的关键突破很大程度来自用强 LM 根据 image captions 和 metadata 生成 conversation、detailed description 与 reasoning questions，再用这些视觉 instruction samples 微调模型；同一套 architecture 能否 OCR、grounding 或做图表推理，往往由任务覆盖决定。

\begin{figure}[H]
\centering
\includegraphics[width=0.82\textwidth]{llava-example.png}
\caption{LLaVA instruction example：图像输入与多轮视觉问答。}
\figsource{Stanford CS336 2026 Lecture 17 官方可执行讲义，\texttt{images/llava-example.png}。}
\end{figure}

\begin{knowledgebox}{Nota de clase}
课程在 LLaVA-OneVision 小结中明确强调，多模态模型的大量工作来自 data curation，而非激进 architecture。Synthetic task-specific data 决定模型是否会 OCR、chart QA、grounding、multi-image reasoning 和 GUI interaction。
\end{knowledgebox}

\subsection{为什么 LLaVA 不能原生生成图像}

回到开头的 understanding/generation 区分，LLaVA 只把图像编码成连续 features 输入 LM，输出仍是 text tokens。若要生成 pixels，通常需要外接 diffusion decoder 或 image generation tool；因此“能看图”不能推出“能在同一 token space 原生画图”。

\begin{figure}[H]
\centering
\includegraphics[width=0.82\textwidth]{llava-gen.png}
\caption{理解与生成的边界：视觉 encoder 能输入图像，但输出图像需要额外 decoder。}
\figsource{Stanford CS336 2026 Lecture 17 官方可执行讲义，\texttt{images/llava-gen.png}。}
\end{figure}

\begin{warningbox}{不要把 tool-mediated generation 当作 native generation}
模型调用 diffusion tool 生成图片，与一个统一 autoregressive model 直接输出 image tokens 是不同系统。前者模块化、质量高；后者统一但训练和离散化更难。
\end{warningbox}

\subsection{本章小结}

LLaVA 建立了现代 VLM 的标准模板：预训练视觉 encoder、轻量 projector、LLM 与视觉 instruction data。它擅长 understanding，但 generation 仍需要外部模型。

